\chapter{Sensitivity of the Frozen Adjoint to Forward State Convergence}
\label{app:picard_sensitivity}

This appendix provides supplementary data for the sensitivity study regarding the number of Picard iterations ($N_{Picard}$) used in the decoupled fluid-turbulence solver under the frozen adjoint formulation, as referenced in Section 4.8. 

\section{Visual Comparison of Topologies}
A direct comparison of the final projected topologies ($\bar{\gamma}$) for all three benchmarks is provided in Figure~\ref{fig:app_visual_all}. In all cases, an under-converged forward state ($N_{Picard} = 1$) resulted in [describe general trend, e.g., artificial blockages, jagged structural boundaries, or less streamlined conduits]. Conversely, enforcing strict forward equilibrium ($N_{Picard} = 3$) successfully recovered aerodynamically refined profiles that align with the physically coupled (semi-frozen) adjoint results.

\begin{figure}[htpb]
	\centering
	% Diffuser Row
	\begin{subfigure}[b]{0.45\textwidth}
		\centering
		\includegraphics[width=\textwidth]{example-image}
		\caption{Diffuser: $N_{Picard}=1$}
	\end{subfigure}
	\hfill
	\begin{subfigure}[b]{0.45\textwidth}
		\centering
		\includegraphics[width=\textwidth]{example-image}
		\caption{Diffuser: $N_{Picard}=3$}
	\end{subfigure}
	
	\vspace{1em}
	
	% Pipe Bend Row
	\begin{subfigure}[b]{0.45\textwidth}
		\centering
		\includegraphics[width=\textwidth]{example-image}
		\caption{Pipe Bend: $N_{Picard}=1$}
	\end{subfigure}
	\hfill
	\begin{subfigure}[b]{0.45\textwidth}
		\centering
		\includegraphics[width=\textwidth]{example-image}
		\caption{Pipe Bend: $N_{Picard}=3$}
	\end{subfigure}
	
	\vspace{1em}
	
	% Double Pipe Row
	\begin{subfigure}[b]{0.45\textwidth}
		\centering
		\includegraphics[width=\textwidth]{example-image}
		\caption{Double Pipe: $N_{Picard}=1$}
	\end{subfigure}
	\hfill
	\begin{subfigure}[b]{0.45\textwidth}
		\centering
		\includegraphics[width=\textwidth]{example-image}
		\caption{Double Pipe: $N_{Picard}=3$}
	\end{subfigure}
	\caption{Impact of Picard iteration count on topological evolution across all benchmark cases.}
	\label{fig:app_visual_all}
\end{figure}

\section{Numerical Performance Metrics}
Table~\ref{tab:app_metrics_all} summarizes the final objective function values ($J$) representing total power dissipation, alongside the computational wall-clock time required per optimization (MMA) iteration. 

\begin{table}[htpb]
	\centering
	\caption{Comparative performance metrics for $N_{Picard}=1$ and $N_{Picard}=3$.}
	\label{tab:app_metrics_all}
	\renewcommand{\arraystretch}{1.2}
	\begin{tabular}{llcc}
		\hline
		\textbf{Benchmark} & \textbf{Metric} & \textbf{$N_{Picard}=1$} & \textbf{$N_{Picard}=3$} \\
		\hline
		\textbf{Diffuser} & Final Objective ($J$) & [Value] & [Value] \\
		& Avg. Time/MMA Iter [s] & [Value] & [Value] \\
		\hline
		\textbf{Pipe Bend} & Final Objective ($J$) & [Value] & [Value] \\
		& Avg. Time/MMA Iter [s] & [Value] & [Value] \\
		\hline
		\textbf{Double Pipe} & Final Objective ($J$) & [Value] & [Value] \\
		& Avg. Time/MMA Iter [s] & [Value] & [Value] \\
		\hline
	\end{tabular}
\end{table}